\chapter{Lifted inference}\label{cha:probabilisticRelationalModels}

\abstract{
    Computation networks in \firstOrderLogic{} contain incoming axes to the substitution and the semantic structure.
    It is often not possible to contract them and therefore to perform reasoning schemes on them.
    We in this chapter investigate methods to reduce reasoning to the activation network, a strategy which we call \emph{lifting}.
    First of all, we review inference schemes developed in \firstOrderLogic{}.
    We further investigate, how \ComputationActivationNetworks{} can be trained given data in the form of a first-order logic world.
}

%\chapter{\chatextfolModels}\label{cha:folModels}
%
%We in this chapter extend the tensor network based treatment of \propositionalLogic{} towards \firstOrderLogic{}.
%In contrast to \propositionalLogic{}, worlds in \firstOrderLogic{} contain objects, which have relations between each others.
%We show in this chapter, how these relations are captured by the \substitutionStructure{} of each world and derive tensor representations of each world.
%The \substitutionStructure{} is then combined with the \semanticStructure{}, which enumerates possible worlds in a theory.
%We then generalize \HybridLogicNetworks{} to the situation of \firstOrderLogic{}s and show in particular, that the likelihood of \HybridLogicNetworks{} on a single \firstOrderLogic{} world is in some cases equivalent to the likelihood of a dataset in propositoinal logics.


\section{Ideas}

We here understand lifting as the reduction of inference in structured representations to factored representations.
%% CAN as a structure for lifting
In the perspective of \ComputationActivationNetworks{}, we find lifting schemes as reductions to inference in the activation network.
We draw on the representations investigated in \charef{cha:firstOrderLogic}, where the computation network contains the FOL atoms, and the activation the connectives acting on them.


We investigate two lifting schemes:
\begin{itemize}
    \item \textbf{FOL lifted inference:}
    Restrict the computation networks implicitly by contraction equations and restrict to the activation network parts of these equations.
    Simplify the activation network based on substitutions and propositional inference.
    \item \textbf{Probabilistic Relational Models (PRM):}
    Represent the likelihood of a single observed FOL world by the likelihood of a dataset in a factored representation.
\end{itemize}


\subimport{inference/}{inference}
\subimport{probabilistic-relational-models/}{probabilistic-relational-models}
